\documentclass[11pt]{article}
\usepackage[margin=0.88in]{geometry}
\usepackage{amsmath,amssymb,amsthm,mathtools,booktabs}

\newcommand{\R}{\mathbb R}
\newcommand{\1}{\mathbf 1}
\newcommand{\Tr}{\operatorname{Tr}}
\newcommand{\norm}[1]{\lVert#1\rVert}
\newcommand{\ip}[2]{\langle #1,#2\rangle}
\newcommand{\diag}{\operatorname{diag}}
\newtheorem{proposition}{Proposition}

\title{Which Optimization Produced the Rank-One Majority Witness?\\
The SDP, the Half-Set Search, and the Numerical $4/\pi$ Law}
\author{}
\date{}

\begin{document}
\maketitle

\begin{abstract}
For odd $k$, this note records the precise optimization path leading to the
period-four rank-one feasible family for Majority.  The convex symmetric SDP
is a trace-minimization problem, equivalently a nuclear-norm problem for the
cross-cut block.  It is useful as a benchmark, but its optimizer is generally
not rank one.  Imposing rank one gives a nonconvex bilinear program.  In the
Krawtchouk basis its $k$ feasibility equations become coordinatewise
complementarity, reducing the search to a finite optimization over half-sets
of Fourier degrees.  Exhaustive and periodic searches select the period-four
half-set.  Exact evaluation of that explicit witness gives
$\mathcal C_k-(4/\pi)\log k\approx2.733$.

The logical qualification is important: the current written proof certifies
$\mathcal C_k=O(\log k)$, while the coefficient $4/\pi$ is a very stable
numerical asymptotic of the explicit rank-one formula.  It was not obtained
as the certified optimum value of a separate SDP.
\end{abstract}

\section{The short answer}

The calculation had the following order:
\[
\begin{gathered}
\boxed{\text{symmetric PSD SDP}}
\ \longrightarrow\
\boxed{\text{rank-one bilinear restriction}}
\ \longrightarrow\
\boxed{\text{Fourier half-set optimization}}\\[2mm]
\longrightarrow\
\boxed{Z_4=\{0,1\}\pmod 4}
\ \longrightarrow\
\boxed{\mathcal C_k\sim\tfrac4\pi\log k\ \text{numerically}}.
\end{gathered}
\]
Thus the optimization that actually produced the rank-one formula was the
half-set problem \eqref{eq:halfset-program} below.  The SDP in
\eqref{eq:count-sdp}, or its nuclear-norm form \eqref{eq:nuclear}, motivated
and benchmarked the search, but did not itself return the period-four
rank-one solution.  The value $4/\pi$ was recognized only after exact costs
of the period-four member were evaluated for large $k$.

\section{From the dual adversary SDP to a $(k+1)$-dimensional SDP}

Let $k=2m+1$.  Concepts are $k$-sets $A$, queries are subsets $S$, and
$f_A(S)$ is the Majority label determined by $|A\cap S|$.  The starting dual
adversary program is
\begin{equation}
 \begin{aligned}
 \text{minimize}\quad & \max_A\sum_S X_S[A,A],\\
 \text{subject to}\quad &X_S\succeq0,\\
 &\sum_{S:f_A(S)\ne f_B(S)}X_S[A,B]=1
       \qquad(A\ne B).
 \end{aligned}                                           \tag{D}\label{eq:dual}
\end{equation}

Put
\[
 \pi_x=2^{-k}\binom{k}{x},\qquad L=\{0,\ldots,m\},
 \qquad R=\{m+1,\ldots,k\}.
\]
If two concepts are at Johnson distance $d=|A\setminus B|$, define their
normalized joint-count kernel
\begin{equation}
 T_d[x,y]=\frac{\Pr\{|A\cap S|=x,\ |B\cap S|=y\}}
                 {\sqrt{\pi_x\pi_y}}.                   \label{eq:kernel}
\end{equation}
Equivalently, if $C,P,Q$ are independent fair binomials of sizes
$k-d,d,d$, then the numerator is $\Pr(C+P=x,C+Q=y)$.

Permutation symmetrization reduces \eqref{eq:dual} to the count-space SDP
\begin{equation}
 \boxed{
 \begin{aligned}
  \text{minimize}\quad &\Tr G,\\
  \text{subject to}\quad&G\in\mathbb S_+^{k+1},\\
  &2\ip{T_d[L,R]}{G[L,R]}_F=1,\qquad d=1,\ldots,k .
 \end{aligned}}                                         \tag{P$_k$}\label{eq:count-sdp}
\end{equation}
This is the concrete SDP implemented in
\texttt{majority\_direct\_count\_sdp.py}.

It has an exact cross-block reformulation.  Write $Q=G[L,R]$.  The minimum
trace of a PSD matrix with prescribed off-diagonal block $Q$ equals
$2\norm{Q}_*$.  Indeed, if $Q=U\Sigma V^{\mathsf T}$, the completion
\[
 \begin{pmatrix}U\Sigma U^{\mathsf T}&Q\\
 Q^{\mathsf T}&V\Sigma V^{\mathsf T}\end{pmatrix}
 =\begin{pmatrix}U\\V\end{pmatrix}\Sigma
  \begin{pmatrix}U\\V\end{pmatrix}^{\!\mathsf T}
\]
is PSD and has trace $2\Tr\Sigma$; the reverse inequality is the standard
PSD block/nuclear-norm inequality.  Hence \eqref{eq:count-sdp} is equivalent
to
\begin{equation}
 \boxed{
 \begin{aligned}
  \text{minimize}\quad&2\norm Q_*,\\
  \text{subject to}\quad&2\ip{T_d[L,R]}Q_F=1,
                   \qquad d=1,\ldots,k .
 \end{aligned}}                                         \tag{N$_k$}\label{eq:nuclear}
\end{equation}

This convex program does not explain rank one by itself.  For example, its
numerical cross-block ranks for $k=3,5,7,9$ are respectively $2,3,3,4$
(at a relative singular-value cutoff of $10^{-7}$).  Rank one therefore has
to be imposed as additional structure.

\section{The actual rank-one optimization}

Write a rank-one PSD variable as $G=zz^{\mathsf T}$ with
$z=(u,v)\in\R^L\oplus\R^R$.  Program \eqref{eq:count-sdp} becomes
\begin{equation}
 \begin{aligned}
 \text{minimize}\quad&\norm{u}^2+\norm{v}^2,\\
 \text{subject to}\quad&2u^{\mathsf T}T_d[L,R]v=1,
                         \qquad d=1,\ldots,k .
 \end{aligned}                                          \tag{R$_k$}\label{eq:rankone}
\end{equation}
This is a nonconvex quadratically constrained problem.  Rescaling
$(u,v)\mapsto(cu,c^{-1}v)$ and minimizing over $c>0$ gives an equivalent
scale-invariant form.  It is convenient to write
\begin{equation}
 u=\sqrt{\frac{\norm b}{2\norm a}}\,a,
 \qquad
 v=\sqrt{\frac{\norm a}{2\norm b}}\,b .               \label{eq:balance}
\end{equation}
Then the program is exactly
\begin{equation}
 \boxed{
 \begin{aligned}
 \text{minimize}\quad&\norm a\,\norm b,\\
 \text{subject to}\quad&a^{\mathsf T}T_d b=1,
                         \qquad d=1,\ldots,k,\\
 &\operatorname{supp}(a)\subseteq L,
 \quad \operatorname{supp}(b)\subseteq R.
 \end{aligned}}                                         \tag{B$_k$}\label{eq:bilinear}
\end{equation}
Here $a,b$ are viewed as $(k+1)$-vectors after zero extension.  Its objective
is precisely $\norm z^2=\norm a\norm b$.

\section{Why the bilinear equations become a half-set search}

Let $U$ be the orthogonal normalized Krawtchouk matrix,
\[
 U[x,t]=2^{-k/2}\sqrt{\frac{\binom{k}{x}}{\binom{k}{t}}}
 \sum_j(-1)^j\binom{x}{j}\binom{k-x}{t-j}.
\]
The kernels diagonalize simultaneously:
\begin{equation}
 T_d=U\diag\!\left(
 \frac{\binom{k-d}{t}}{\binom{k}{t}}
 \right)_{t=0}^kU^{\mathsf T}.                          \label{eq:diagonalization}
\end{equation}
Set $\alpha=U^{\mathsf T}a$ and $\beta=U^{\mathsf T}b$.  Because the
functions $d\mapsto\binom{k-d}{t}/\binom{k}{t}$ for $0\le t<k$ are
triangularly independent, all $k$ equations in \eqref{eq:bilinear} are
equivalent to
\begin{equation}
 \alpha_0\beta_0=1,\qquad
 \alpha_t\beta_t=0\quad(1\le t\le k-1).               \label{eq:complementarity}
\end{equation}
The $t=k$ multiplier vanishes for every $d\ge1$.  Moreover
$a^{\mathsf T}b=0$ because the count supports are disjoint, so
$\alpha_k\beta_k=-1$ follows automatically.

Normalize $\alpha_0=\beta_0=1$.  In a nondegenerate solution each interior
degree must therefore be assigned to one of two zero sets.  Choose an
$m$-element set $Z\subset I:=\{1,\ldots,k-1\}$ and let $Z^c=I\setminus Z$.
Define the square Krawtchouk submatrices
\[
 B_Z=U[L,\{0\}\cup Z],\qquad
 C_Z=U[R,\{0\}\cup Z^c]
\]
and
\begin{equation}
 a_L=B_Z^{-\mathsf T}e_0,\quad a_R=0,\qquad
 b_R=C_Z^{-\mathsf T}e_0,\quad b_L=0.                  \label{eq:ab}
\end{equation}
Then $\alpha_0=\beta_0=1$, $\alpha_t=0$ on $Z$, and $\beta_t=0$ on
$Z^c$, so \eqref{eq:complementarity} holds.  The continuous rank-one search
has consequently reduced to the finite optimization
\begin{equation}
 \boxed{
 \min_{\substack{Z\subset\{1,\ldots,k-1\}\\|Z|=m}}
 \mathcal C(Z),\qquad
 \mathcal C(Z)=\norm{B_Z^{-\mathsf T}e_0}\,
                \norm{C_Z^{-\mathsf T}e_0}.}           \tag{H$_k$}\label{eq:halfset-program}
\end{equation}
This is the optimization implemented with exact rational interpolation
weights in the accompanying script
\[
 \texttt{majority\_rank\_one\_search.py}.
\]

\section{How the period-four formula emerged}

Exhaustive evaluation of \eqref{eq:halfset-program} for every odd
$k\le15$ selected, up to reflection and complementation,
\begin{equation}
 \boxed{Z_4=\{t:1\le t\le k-1,\ t\equiv0\text{ or }1\pmod4\}.}             \label{eq:period4}
\end{equation}
Searches over balanced periodic masks of periods $6$ and $10$ through
$k=63$ did not improve it.  Thus the simple residue-class description was a
pattern discovered in the exact half-set optimization; it was then promoted
to a formula valid for every odd $k$.

For completeness, the resulting query matrices have the one-vector form
\begin{equation}
 \boxed{X_S=2^{-n}(M_Sz)(M_Sz)^{\mathsf T},\qquad
 M_S[A,x]=\frac{\1\{|A\cap S|=x\}}{\sqrt{\pi_x}},}      \label{eq:XS}
\end{equation}
where $a,b$ are given by \eqref{eq:ab} with $Z=Z_4$, and $z=(u,v)$ is
given by \eqref{eq:balance}.  Hence every $X_S$ is visibly PSD and rank one.
Equations \eqref{eq:diagonalization}--\eqref{eq:complementarity} prove all
feasibility constraints, and the objective is $\mathcal C(Z_4)$.

The inverses in \eqref{eq:ab} can be removed entirely.  Cauchy interpolation
expresses the nonzero Fourier entries as finite products over the two
step-four progressions.  Those are the explicit product formulas used by the
exact-arithmetic evaluator and in the written $O(\log k)$ proof.

\section{Origin of the $4/\pi$ coefficient}

The exact period-four cost was evaluated as a rational number squared, with
only the final square root rounded.  For $k\equiv3\pmod4$ the results are:
\begin{center}
\begin{tabular}{rcc}
\toprule
$k$&$\mathcal C_k$&$\mathcal C_k-(4/\pi)\log k$\\
\midrule
31   & 7.094145520 & 2.721857215\\
63   & 8.002838568 & 2.727635595\\
127  & 8.898138883 & 2.730328322\\
255  & 9.786974808 & 2.731618934\\
511  &10.672640712 & 2.732248334\\
1023 &11.556738512 & 2.732558512\\
2047 &12.440056881 & 2.732712325\\
4095 &13.322986795 & 2.732788875\\
\bottomrule
\end{tabular}
\end{center}
A least-squares fit of $\mathcal C_k=A\log k+B$ on these eight points gives
\begin{equation}
 A=1.27510935,\qquad B=2.71921682,\qquad
 \frac4\pi=1.27323954.                                  \label{eq:fit}
\end{equation}
More tellingly, the third column stabilizes rapidly.  This led to the
empirical law
\begin{equation}
 \boxed{\mathcal C_k=\frac4\pi\log k+c+o(1),
 \qquad c\approx2.7329.}                                \label{eq:conjecture}
\end{equation}

The status of \eqref{eq:conjecture} should not be overstated.  The finite
products for the period-four vector, combined with Wallis and Gautschi
bounds, currently give the rigorous estimate
\begin{equation}
 \mathcal C_k\le16H_{m+1}=O(\log k).                    \label{eq:proved}
\end{equation}
They do not yet give the sharp leading constant $4/\pi$.  A proof of that
constant would require a sharp asymptotic evaluation of the four
Gamma-product residue classes, with a summable error uniform from the bulk
to both endpoints.  In particular, the honest provenance is:
\[
 \text{half-set search}\Rightarrow Z_4
 \Rightarrow \text{exact cost data}\Rightarrow 4/\pi\text{ conjecture},
\]
not ``an SDP solver certified the coefficient $4/\pi$.''

\section{Reproduction}

From the \texttt{majority} directory, the relevant computations are
\begin{verbatim}
python3 majority_direct_count_sdp.py --ks 3 5 7 9
python3 majority_rank_one_search.py --ks 3 5 7 9 11 13 15 --exhaustive
python3 majority_rank_one_search.py --ks 31 63 127 255 511 1023 2047 4095
\end{verbatim}
The first command solves the convex SDP/nuclear-norm benchmark.  The second
performs the half-set search that identifies the period-four mask.  The third
evaluates its explicit exact-arithmetic formula at large $k$; regressing those
values produces \eqref{eq:fit}.

\end{document}
